\chapter{Введение}
\label{ch:intro}

В промышленных процессах, связанных с транспортировкой и хранением газов, значительные энергозатраты возникают из-за неконтролируемого изменения температуры газа при дросселировании. Это явление особенно критично в системах газораспределения и криогенных установках, где даже небольшие колебания температуры могут привести к снижению эффективности работы оборудования или повреждению инфраструктуры. Проблема усугубляется при работе с реальными газами, поведение которых существенно отличается от идеальных моделей, что затрудняет прогнозирование температурных эффектов.

Первым способом решения данной проблемы может быть использование систем активного охлаждения или подогрева газа, которые будут компенсировать нежелательные температурные изменения. Однако такой подход требует значительных энергетических затрат и сложного оборудования, что увеличивает эксплуатационные расходы и снижает надежность системы.

Иначе подойти к решению этой задачи позволяет знание коэффициента Джоуля-Томсона. Коэффициент Джоуля-Томсона характеризует изменение температуры газа при адиабатическом дросселировании. Это дает возможность точно прогнозировать температурные эффекты на основе термодинамических свойств рабочего тела, что позволяет оптимизировать параметры технологических процессов без дополнительных энергозатрат, выбирая оптимальные давления и температуры дросселирования. Такой подход не только повышает энергоэффективность установок, но и обеспечивает более безопасную эксплуатацию за счет предотвращения нежелательных температурных скачков.

Целью работы являлась проверка метода определения коэффициента Джоуля-Томсона и параметров адиабатического расширения для углекислого газа с применением модели газа Ван-дер-Ваальса.